\subsection{N-024: exact minima in a specified Laurent-parameter family}

The all-table minimum in the user's question remains open.  We prove
an all-degree minimum in the precise two-parameter family below:
3898 pairs without bar invariance, and4046 with bar invariance.
Both bounds are attained by more than one table.  These are also new
upper bounds for the corresponding unrestricted problems.  No
unimodality or preferred categorical realization is asserted.
The final extension below proves the same minima when the common
cell scale itself is allowed to be an arbitrary Laurent polynomial;
distinct cell scales and the other fixed parameters are not varied.

\paragraph{The family, including its unrestricted parameters.}
Write $A=\mathbb Z[q,q^{-1}]$, $D=q+2+q^{-1}$,
$u=D-2=q+q^{-1}$ and $a=(D-1)u$.
Keep the ten labels $\mu,\Sigma,\gamma^i_x$, where $i=1,2$ and
$x\in\{1,12,2,21\}$.  The parameters $W,L$ range over all of $A$;
neither parameter is required to be coefficientwise nonnegative.
Put
\[
 K_1=u,\quad K_2=D,\quad G_1=2Wa,\quad G_2=Wa,\quad S=2Wa^2.
 \tag{N-151}
\]
Give $\mu$ its strict unit products and use
\[
\begin{aligned}
 \Sigma^2&=S\Sigma,&
 \Sigma\gamma^i_x=\gamma^i_x\Sigma&=K_i^{|x|-1}G_i\Sigma,\\
 \gamma^i_x\gamma^j_y&=W K_i^{|x|-1}K_j^{|y|-1}\Sigma
       &&(i\ne j).
\end{aligned}
 \tag{N-152}
\]
For a same-cell product, reduce the concatenated word $xy$ using
the two relations of lengths two and three.  Let $\alpha$ count
adjacent-letter deletions, $\beta$ count alternating-triple reductions,
and $z\in\{1,12,2,21\}$ be the reduced word.  Here
$\alpha,\beta\in\{0,1\}$ and
$|x|+|y|=|z|+\alpha+2\beta$.
The four types $(0,0),(1,0),(0,1),(1,1)$ occur with multiplicities
$2,6,6,2$.  Define
\[
\begin{aligned}
 \gamma^1_x\gamma^1_y
   &=aD^\alpha u^\beta\gamma^1_z+
       \bigl(2Wu^{|x|+|y|-2}-D^\alpha u^{\beta+|z|-1}\bigr)\Sigma,\\
 \gamma^2_x\gamma^2_y
   &=aD^\alpha L^\beta\gamma^2_z+
    D^{|z|+\alpha-2+\beta}
       \begin{cases}(W-D)/2,&\beta=0,\\(DW-L)/2,&\beta=1\end{cases}\Sigma.
\end{aligned}\tag{N-153}
\]
The exponent preceding the second bottom coefficient is zero or one.
This specifies every coefficient, including every zero coefficient;
no other dimensions or multiplication parameters vary in the family.

An \emph{admissible table} in this family has integral Laurent
coefficients $d$ satisfying $d\ge|c|$ and $d\equiv c\pmod2$
coefficientwise for the supplied 132-entry trace table $c$.
The absolute value here is coefficientwise.  A bar-invariant table
additionally has each coefficient fixed by $q\mapsto q^{-1}$,
with the ten labels fixed.

\paragraph{Associativity for every specialization.}
Over $F=\mathbb Q(q,W,L)$ use the matrices
\[
 A_i=\begin{pmatrix}D&1\\0&0\end{pmatrix},\qquad
 B_i=\begin{pmatrix}0&0\\\Lambda_i&D\end{pmatrix},
 \qquad\Lambda_1=u,\quad\Lambda_2=L.
 \tag{N-154}
\]
They have quadratic parameter $D$ and triple parameter $\Lambda_i$.
The four word matrices $A_i,A_iB_i,B_i,B_iA_i$ are independent
over this field, because their determinant is nonzero when
$\Lambda_i(D^2-\Lambda_i)\ne0$.  Map the ten labels to
$F_{\rm aug}\times F_{\rm char}\times M_2(F)^{\oplus2}$:
$\mu$ is its unit, $\Sigma=(0,S,0,0)$, and $\gamma^i_x$ has
character component $K_i^{|x|-1}G_i$ and matrix component
$aP_x$ in factor $i$, with every other component zero.
This is a vector-space isomorphism.  Successively the augmentation,
matrix factors and character factor detect its ten coordinates.

The character residual of a same-cell product is
\[
 \frac{K_i^{|x|+|y|-2}G_i^2
       -aD^\alpha\Lambda_i^\beta K_i^{|z|-1}G_i}{S}.
 \tag{N-155}
\]
For $i=1$ this is the first bottom coefficient of (N-153).
For $i=2$, the length identity gives its displayed second coefficient.
The cross residual is (N-152) since $G_1G_2/S=W$.
Thus the table is associative as an identity over
$\mathbb Q[q,q^{-1},W,L]$, and consequently after every Laurent
specialization of $W,L$, including any specialization at which this
particular matrix realization degenerates.  On the admissible
parameters its constants lie in $A$, so it defines an associative
$A$-algebra on the specified free ten-element module.  Word reversal
is an anti-involution and the simultaneous letter interchange
$1\leftrightarrow2$ is an automorphism, for all parameters.

\paragraph{Exact coefficientwise description of admissibility.}
Set
\[
\begin{aligned}
 \delta&=(D-1)(D-4),&
 K&=2D^2-5D+4,\\
 R&=3D^2+2D,&
 H_F&=-D^2+10D-8+4DF,\\
 T&=|\delta u^2|,\\
 T&=q^4+q^{-4}+q^3+q^{-3}
       +2(q^2+q^{-2})+3(q+q^{-1})+2.
\end{aligned}\tag{N-156}
\]
The table is admissible if and only if
\[
 \boxed{\begin{gathered}
 W=3D+4F,\quad F\in\mathbb N[q,q^{-1}],\\
 L\le H_F,\quad L\equiv R\pmod4,\quad aL\ge T.
 \end{gathered}}
 \tag{N-157}
\]
There is no support or degree bound in this assertion.

For necessity, the actual second-cell bottom traces with no length
factor are $D$ and $K$, and the respective dimension coefficients
are $U_2=(W-D)/2$ and $V_2=(DW-L)/2$.
Thus $U_2\ge D$ and $U_2\equiv D\pmod2$ give
$U_2=D+2F$ for a nonnegative integral Laurent $F$, equivalently
$W=3D+4F$.  Likewise $V_2\ge K$ and $V_2\equiv K\pmod2$
give $L\le H_F$ and $L\equiv H_F\equiv R\pmod4$.
The second-cell triple coefficient has trace of coefficientwise
absolute value $T$ and dimension $aL$, proving the last inequality.

For sufficiency, these same conditions give integral admissible
$U_2,V_2$, and their multiples by the nonnegative polynomial $D$.
They also give the second-cell coefficients $aL,aDL$:
$aDL\ge DT\ge|D\delta u^2|$.
Their parity holds since $a\equiv\delta\pmod2$ and
$L\equiv D^2\equiv u^2\pmod2$.
The remaining cell coefficients are $a,aD,au,aDu$;
their four comparisons are exactly the first four explicitly proved
Laurent comparisons in (N-132), with the middle two interchanged.

For completeness the remaining baseline comparisons are displayed
below at $W=3D$.  Each row lists all Laurent coefficients of
$(d-t)/2$, from the negative extreme exponent to the positive one;
the lists are centered at zero.  Every $t$ in this table is already
coefficientwise nonnegative.  Put
$U=2W-D$, $U_1=2W-u$, $J=u(2W-1)$, and $V=2Wu-D$.
\[
\begin{array}{c|l}
 (d,t)&\text{coefficients of }(d-t)/2\\\hline
 (W,D)&(1,2,1)\\
 (uW,D^2)&(1,1,0,1,1)\\
 (G_1,4D(D-1))&(3,7,9,10,9,7,3)\\
 (G_2,D^2(D-1))&(1,2,2,2,2,2,1)\\
 (uG_1,4D^2(D-1))&(3,7,8,5,2,5,8,7,3)\\
 (S,4D^2(D-1)^2)&(3,10,18,20,15,12,15,20,18,10,3)\\
 (U,D)&(2,4,2)\\
 (U_1,D)&(2,5,2)\\
 (uU,D^2)&(2,3,2,3,2)\\
 (J,5D-4)&(3,3,3,3,3)\\
 (uJ,D(5D-4))&(3,3,1,0,1,3,3)\\
 (uV,D(5D-4))&(3,3,0,0,0,3,3).
\end{array}\tag{N-158}
\]
The six first-cell bottom types in the last six rows occur with
multiplicities $2,2,4,4,2,2$, as in N-022.  All unlisted cross
or scalar length variants follow by multiplication by $D$.
Replacing $W=3D$ by $3D+4F$ increases each of these dimension
coefficients by twice a nonnegative integral Laurent polynomial.
This proves every remaining domination and parity condition and
completes the equivalence (N-157).

\paragraph{Lower bounds with arbitrary Laurent support.}
Let $f=F(1)$ and $\ell=L(1)$ for any admissible parameters.
We claim
\[
             f\ge2,\qquad \ell\ge4,
       \qquad \ell\ge8\text{ if the table is bar invariant}.
             \tag{N-159}
\]
First, $aL\ge T\ne0$ implies that the lowest and highest nonzero
Laurent coefficients of $L$ are positive: the extreme coefficients
of $a$ are one, and the product $aL$ is nonnegative.
If $f=0$, then $F=0$ and $H_F$ has coefficient $-1$ at each
of $q^{-2},q^2$, and vanishes beyond these exponents.  The inequality
$L\le H_F$ is incompatible with those positive extreme coefficients.
If $f=1$, nonnegativity and integrality force $F=q^N$ for some integer
$N$.  When $N\ge0$, the term $4DF$ has no exponent at most $-2$.
Hence $H_F$ has coefficient $-1$ at $-2$ and vanishes below $-2$.
The inequality forces $L_{-2}\le-1$, so the lowest nonzero
coefficient of $L$ has exponent at most $-2$, where it cannot be
positive.  This is a contradiction.  When $N\le0$, the identical
argument at the upper extreme $2$ applies.  Thus $f\ge2$ for every
support, without a bounded search.

Next, $a(1)=6$ and $T(1)=16$ give $6\ell\ge16$.
The congruence in (N-157) gives $\ell\equiv R(1)=56\equiv0\pmod4$,
so $\ell\ge4$.

For the stronger bound suppose the full table is bar invariant.
Because $W$ and $aL$ occur as actual structure constants, $W$ and
$aL$ are bar invariant.  Since $a$ is bar invariant and nonzero,
$L$ is bar invariant too.  Consider the least coefficientwise
nonnegative Laurent polynomial above $T$ and congruent to $aR$
modulo four.  Directly at its nine supported exponents it is
\[
 C=T+2(q^4+q^{-4}),\qquad
 (C_{-4},\ldots,C_4)=(3,1,2,3,2,3,2,1,3),\quad C(1)=20.
 \tag{N-160}
\]
At other exponents both the lower bound and the prescribed residue
are zero.  Consequently every admissible $aL$ has the form
$aL=C+4E$ with $E\in\mathbb N[q,q^{-1}]$.
If $\ell=4$, then $E(1)=(24-20)/4=1$, so $E=q^N$.
At $q=i$ one has $a(i)=0$, $C(i)=4$, whence $i^N=-1$.
At a primitive cube root $\omega$, one has $a(\omega)=0$,
$C(\omega)=-4$, whence $\omega^N=1$.
Thus $N\equiv6\pmod{12}$.  On the other hand bar invariance of
$aL$ and $C$ forces $q^N=q^{-N}$ in $A$, so $N=0$, a contradiction.
Since $\ell$ is a positive multiple of four, $\ell\ge8$.
Only the fourth-root equation is needed for this contradiction;
the cube-root equation further identifies the possible nonbar residue.

\paragraph{Score and exact attainment.}
At $q=1$ write $w=W(1)=12+4f$.  The seven nonunit dimension-mass
groups are
\[
\begin{array}{c|r}
\Sigma^2&72w\\
\Sigma\gamma\text{ and }\gamma\Sigma&264w\\
\text{cross-cell products}&120w\\
\text{first-cell bottom coefficients}&72w-76\\
\text{second-cell bottom coefficients}&50w-40-10\ell\\
\text{first-cell nonbottom coefficients}&324\\
\text{second-cell nonbottom coefficients}&156+84\ell.
\end{array}\tag{N-161}
\]
For example the second bottom mass is
$4(1+4)((w-4)/2+(4w-\ell)/2)$, and the second cell mass is
$6(2+6\cdot4+6\ell+2\cdot4\ell)$.
The first bottom mass follows from the six multiplicities following
(N-158).  Their sum is $578w+74\ell+364$.
The fixed nonunit absolute trace mass is4424, independently read
from all132 supplied expressions.  The19 unit entries contribute
zero pairs.  Thus the total score is
\[
 \operatorname{pairs}(d,c)
    =\frac{578w+74\ell+364-4424}{2}
    =1438+1156f+37\ell.
 \tag{N-162}
\]
The bounds (N-159) imply scores at least3898 and4046, respectively.

For the general minimum take
\[
\begin{aligned}
 F&=q^{-1}+q^3,\qquad W=3D+4F,\\
 L&=4q^4-4q^3-q^2+6q-2-2q^{-1}+3q^{-2}.
\end{aligned}\tag{N-163}
\]
Then $f=2$, $\ell=4$, and the following exact identities prove
all the conditions in (N-157):
\[
 L\equiv R\pmod4,\qquad aL=C+4q^6,\qquad
 H_F-L=4(4q^{-1}+3+q^2+3q^3)\ge0.
\]
This table has3898 pairs, including3654 in the bottom output.
Replacing $q$ by $q^{-1}$ gives a distinct table attaining the
same minimum.  The negative coefficients of $L$ are permitted:
$L$ itself is an auxiliary matrix parameter, while the actual
cell structure constants are $aL$ and $aDL$, both nonnegative.

For bar-invariant minima take $F=u$ and $W=7D-8$, with either
\[
 \begin{aligned}
 L_A&=3D^2-10D,& H_F-L_A&=12D-8,\\
 L_B&=3D^2-14D+16,& H_F-L_B&=16D-24.
 \end{aligned}\tag{N-164}
\]
Both right-hand differences are coefficientwise nonnegative and
divisible by four, and $L_A,L_B\equiv R\pmod4$.
The complete centered coefficient lists of their cell products are
\[
 aL_A:(3,5,6,7,6,7,6,5,3),\qquad
 aL_B:(3,1,10,3,14,3,10,1,3).
\]
Both dominate $T$.  Thus these two distinct bar-invariant tables
have $f=2$, $\ell=8$ and4046 pairs each, with3634 in the bottom
output.  All four examples preserve the Dynkin interchange and
word-reversal symmetries.  This proves the two claimed minima and
their nonuniqueness within exactly the specified family.

\paragraph{A smaller cell scale cannot improve these bounds in this ansatz.}
One nearby attempt replaces $a$ everywhere in (N-151)--(N-155) by
$a_0=|\delta|=D(D-3)$, leaving every other parameter choice and
formula unchanged.  Every admissible table in this modified family
has at least4282 pairs, even without bar invariance.  This is a
lower bound, not a claim of attainment in the modified family.

Indeed the unchanged second-cell bottom conditions still force
$W=3D+4F$ with $F\ge0$ integral.  Put $f=F(1)$ and
$r=F_{-2}+F_{-1}+F_1+F_2$.
The coefficients of $a_0$ at $\pm1,\pm2$ are one and all its
other coefficients vanish.  The true second scalar coefficient
$D^2(D-1)$ has coefficient14 at zero.  Its proposed value is
$G_2=Wa_0$, whose coefficient there is $6+4r$.
Consequently
\[
                              r\ge2.
                              \tag{N-165}
\]
The centered coefficients of $a_0^2$ are
$(1,2,1,2,4,2,1,2,1)$.  Its central coefficient is4, and every
other coefficient is at most2.  The true $\Sigma^2$ coefficient
$4D^2(D-1)^2$ has central coefficient144.  The proposed coefficient
$S=2Wa_0^2$ satisfies
\[
\begin{aligned}
 [q^0]S
 &=72+8[q^0](Fa_0^2)\\
 &\le72+8\bigl(4(f-r)+2r\bigr)
 \le40+32f.
\end{aligned}\tag{N-166}
\]
The first inequality remains valid even if $F$ has arbitrarily
distant support: each of the $r$ units has convolution weight at
most2, and each remaining unit has weight at most4.
Thus $144\le40+32f$ forces the integer $f\ge4$.
The actual second-cell triple coefficient $a_0L\ge T$ gives
$4\ell\ge16$, so $\ell=L(1)\ge4$.

The same seven mass groups, now with $a_0(1)=4$, sum to
$450w+46\ell+204$, where $w=12+4f$.
They are respectively $32w$, $176w$, $120w$, $72w-76$,
$50w-40-10\ell$, $216$, and $104+56\ell$.
The pair score is therefore
\[
 \frac{450w+46\ell+204-4424}{2}
       =590+900f+23\ell\ge4282.
 \tag{N-167}
\]
This rules out this particular smaller-scale route to improving
either of the attained bounds above.  It imposes no lower bound
on dimension rings outside the modified ansatz.
The short exact check \nolinkurl{dimension_a0_obstruction.py}
verifies these polynomial coefficients directly against the two
actual supplied trace expressions and verifies the mass identity.

\paragraph{The same exact minima with arbitrary common Laurent cell scale.}
Enlarge the family once more: in all of (N-151)--(N-155), replace
every occurrence of $a$ by a single arbitrary
$b\in\mathbb Z[q,q^{-1}]$.  Keep the two quadratic parameters
$D$, the first triple parameter $u$, the character-length parameters
$u,D$, and the relations $G_1=2Wb$, $G_2=Wb$, $S=2Wb^2$.
Only $b,W,L$ now vary.  Then
\[
\begin{array}{c|c}
 \text{admissible tables in this enlarged family}&
                  \text{minimum total pairs}\\\hline
 \text{without bar invariance}&3898\\
 \text{with bar invariance on every coefficient}&4046.
\end{array}\tag{N-168}
\]
Both minima remain nonunique, by the same four witnesses already
constructed.  This statement permits any finite Laurent support
for each of its three parameters.

To prove it, note first that $b$ itself occurs as a nonbottom
structure coefficient whose trace has absolute value $a_0=|\delta|$.
Admissibility therefore forces
\[
 b=a_0+2E,\qquad E\in\mathbb N[q,q^{-1}],\qquad
 m=b(1)=4+2E(1).
 \tag{N-169}
\]
In particular $m$ is an even integer at least4.
The second bottom coefficients remain $(W-D)/2$ and $(DW-L)/2$:
the common $b$ cancels out of the character residual (N-155).
Their domination and parity still force $W=3D+4F$, $F\ge0$,
$L\le H_F$ and $L\equiv R\pmod4$.
The second-cell triple coefficient requires $bL\ge T\ne0$.
Since $b$ is nonnegative with strictly positive extreme nonzero
coefficients, the extreme nonzero coefficients of $L$ are positive.
The proof of $F(1)\ge2$ in (N-159) consequently applies verbatim:
it uses only this extreme positivity and $L\le H_F$, not the
particular coefficients or degree of the original $a$.
Writing $f=F(1)$, $\ell=L(1)$, we therefore have
$w=W(1)=12+4f\ge20$.
Evaluation of $bL\ge T$ gives $m\ell\ge16$.
Together with $\ell\equiv0\pmod4$, this implies $\ell\ge4$.

The seven nonunit mass groups now sum to
$2(m+11)^2w+(14m-10)\ell+80m-116$.
Thus every admissible enlarged-family table has score
\[
 B=(m+11)^2w+(7m-5)\ell+40m-2270.
 \tag{N-170}
\]
For $m\ge6$, all variable coefficients here are positive.
Using $w\ge20$ and $\ell\ge4$ gives
\[
 B\ge20(m+11)^2+68m-2290\ge3898,
\]
with the last expression strictly increasing for $m\ge6$.
When $m=4$, (N-169) forces $E=0$, so $b=a_0$ and (N-167)
gives the stronger lower bound4282.  The general minimum is
therefore3898, attained in (N-163).

For a bar-invariant table, the actual coefficient $b$ is bar fixed.
If $m=6$, then $E(1)=1$ in (N-169), so $E=q^N$.
Since $a_0$ is bar fixed, bar invariance forces $N=0$.
Hence $b=a_0+2=(D-1)(D-2)=a$, and the exact earlier family
bound4046 applies.  If $m=4$, its lower bound4282 is larger.
If $m\ge8$, the general bound just obtained is already at least
$20\cdot19^2+68\cdot8-2290=5474$.
The bar minimum is consequently4046, attained twice in (N-164).

Associativity imposes no hidden restriction in this enlargement:
the same product-algebra proof is a polynomial identity with the
common scale $b$ also indeterminate, and hence specializes to every
admissible choice.  Conversely, only necessary coefficient
conditions were used for the lower bounds; no sufficiency assertion
analogous to (N-157) for arbitrary $b$ was presumed.
The enlargement still fixes a common scale for both cells and all
the other displayed relations.  It is not a classification of
arbitrary dimension rings or of arbitrary distinct-cell-scale
families.

\paragraph{Scope and certificate.}
The effective sign--Tate coefficients $(d+c)/2+\varepsilon(d-c)/2$
compose by the equal-parity-pair argument of N-005 and N-022;
the supplied counting consult, Theorem5, supplies the same abstract
regular-matrix realization.  None of this gives a functor into the
ordinary Hecke category or the derived/Koszul comparison requested
in item9(d).  N-023 excludes the stated effective ordinary comparison;
its positivity hypothesis does not apply automatically to derived
Euler classes.  The minimum over all ten-label associative tables,
and canonical selection at ranks three and four, remain open.

N-021 applies to the new general feasible budget $B=3898$ and the
unchanged true exponent bound $d_0=5$.  Thus some global minimizer
over all integral Laurent tables, without a bar requirement, lies
in the explicit exponent box
\[
              |e|\le23388\cdot4^{3898}.
\]
This integer has7811 binary digits, since
$2^{14}\le23388<2^{15}$ and $4^{3898}=2^{7796}$.
This is a smaller finite-search box, not a computation of that
global minimizer.  The separate symmetric-unimodal bound9327 from
N-021, using the9040-pair example, is unchanged.

The constructor \nolinkurl{dimension_scale_family.py} saves all four
complete tables in \nolinkurl{dimension_scale_family_certificate.json}.
The separate verifier \nolinkurl{verify_dimension_scale_family.py}
imports no constructor or supplied table-building routine.  With
integer Laurent-dictionary arithmetic it directly parses all132
true coefficient expressions and checks, for each witness, all1000
associativity triples, all100 products in (N-154)--(N-155), the strict
unit, both label symmetries, every domination and parity inequality,
the stated bar property, and all scores.  Generic independence is
also checked by a nonzero rational determinant at $q=2$.
It verifies the finite polynomial identities in the lower-bound
proof; the arbitrary-support extreme-coefficient argument is proved
above and is not inferred from computation.  Its receipt is
\nolinkurl{dimension_scale_family_verified.json}.
